\documentclass[11pt, a4paper]{article}

\usepackage[T1]{fontenc}
\usepackage[utf8]{inputenc}
\usepackage{tgpagella}
\usepackage{tgcursor}
\usepackage{microtype}
\usepackage[spanish, es-noshorthands]{babel}
\addto\captionsspanish{\renewcommand{\tablename}{Tabla}}

\usepackage{amsmath, amssymb}
\usepackage[table, xcdraw]{xcolor}
\usepackage{booktabs}
\usepackage[most]{tcolorbox}
\usepackage[margin=2.5cm]{geometry}
\usepackage{fancyhdr}
\usepackage{listings}
\usepackage{tikz}
\usetikzlibrary{shapes.geometric, arrows.meta, positioning, calc}

\lstset{
    language=Python,
    basicstyle=\ttfamily\small,
    backgroundcolor=\color{gray!5},
    frame=single,
    rulecolor=\color{gray!40},
    keywordstyle=\color{blue}\bfseries,
    stringstyle=\color{red!70!black},
    commentstyle=\color{green!60!black}
}

\pagestyle{fancy}
\fancyhf{}
\setlength{\headheight}{16pt}
\lhead{\textbf{UCB Sede Tarija} | Ing. Mecatrónica}
\chead{Robótica (IMT-342) - Solucionario}
\rhead{Semana 6 - Sesión 1}
\lfoot{Prof: M.Sc. Ing. Bernardo Quiroga Turdera}
\rfoot{Página \thepage}
\renewcommand{\headrulewidth}{0.4pt}
\renewcommand{\footrulewidth}{0.4pt}

\definecolor{ucbblue}{RGB}{0, 51, 102}

\begin{document}

\begin{center}
    \Large\textbf{\textcolor{red!80!black}{SOLUCIONARIO Y GUÍA DE CALIFICACIÓN (EC1)}}\\[0.2cm]
    \normalsize Documento Confidencial Docente[cite: 7]
\end{center}

\section*{Problema 1 (20 pts)[cite: 7]}
\textbf{a) [4 pts]} Indica la pose del marco 1 respecto al marco 0. Su bloque rotacional es un giro de $+90^\circ$ sobre $z_0$, y su origen está trasladado en $[2,1,0]^T$. (\textit{2 pts interpretación, 1 pt orientación, 1 pt posición}). \\
\textbf{b) [6 pts]} ${}^{0}T_2 = {}^{0}T_1 {}^{1}T_2$.
\begin{equation*}
    {}^{0}T_2 = \begin{bmatrix} 0 & -1 & 0 & 2 \\ 1 & 0 & 0 & 2 \\ 0 & 0 & 1 & 0 \\ 0 & 0 & 0 & 1 \end{bmatrix}
\end{equation*}
(\textit{2 pts cadena, 2 pts orientación, 2 pts posición final}). \\
\textbf{c) [6 pts]} ${}^{1}T_0 = ({}^{0}T_1)^{-1}$. $R^T = \begin{bmatrix} 0 & 1 & 0 \\ -1 & 0 & 0 \\ 0 & 0 & 1 \end{bmatrix}$. $-R^T p = -\begin{bmatrix} 0 & 1 & 0 \\ -1 & 0 & 0 \\ 0 & 0 & 1 \end{bmatrix}\begin{bmatrix} 2 \\ 1 \\ 0 \end{bmatrix} = \begin{bmatrix} -1 \\ 2 \\ 0 \end{bmatrix}$. 
\begin{equation*}
    {}^{1}T_0 = \begin{bmatrix} 0 & 1 & 0 & -1 \\ -1 & 0 & 0 & 2 \\ 0 & 0 & 1 & 0 \\ 0 & 0 & 0 & 1 \end{bmatrix}
\end{equation*}
(\textit{2 pts $R^T$, 2 pts $-R^Tp$, 2 pts resultado}). \\
\textbf{d) [4 pts]} La secuencia debe seguir la convención semántica $0 \to 1 \to 2$. Multiplicar ${}^1T_2{}^0T_1$ es inconsistente con los marcos y la multiplicación no es conmutativa. (\textit{2 pts semántica, 1 pt no conmutativa, 1 pt orden correcto}).

\section*{Problema 2 (25 pts)[cite: 7]}
\textbf{a) [6 pts]} $\theta_i$: rotación en $z_{i-1}$. $d_i$: traslación en $z_{i-1}$. $a_i$: traslación en $x_i$. $\alpha_i$: rotación en $x_i$. (\textit{Puntaje parcial por asociaciones de ejes correctas}).\\
\textbf{b) [4 pts]} Al ser revoluta, la variable es $q_i = \theta_i$.\\
\textbf{c) [11 pts]} Reemplazando $a=0, d=L, \alpha=-\pi/2 \implies \cos\alpha=0, \sin\alpha=-1$:
\begin{equation*}
    {}^{i-1}T_i = \begin{bmatrix} \cos q_i & 0 & -\sin q_i & 0 \\ \sin q_i & 0 & \cos q_i & 0 \\ 0 & -1 & 0 & L \\ 0 & 0 & 0 & 1 \end{bmatrix}
\end{equation*}
(\textit{3 pts base DH, 2 pts sust $\alpha$, 6 pts por filas correctas}).\\
\textbf{d) [4 pts]} Para $q_i=0 \implies \cos(0)=1, \sin(0)=0$.
\begin{equation*}
    {}^{i-1}T_i(0) = \begin{bmatrix} 1 & 0 & 0 & 0 \\ 0 & 0 & 1 & 0 \\ 0 & -1 & 0 & L \\ 0 & 0 & 0 & 1 \end{bmatrix}
\end{equation*}

\section*{Problema 3 (45 pts)[cite: 7]}
\textbf{a) [15 pts]} (\textit{5 pts c/u: 2 pts trig, 2 pts rotación, 1 pt traslación})
\begin{equation*}
    {}^{0}T_1 = \begin{bmatrix} 1 & 0 & 0 & 0 \\ 0 & 0 & -1 & 0 \\ 0 & 1 & 0 & 1 \\ 0 & 0 & 0 & 1 \end{bmatrix}, \quad
    {}^{1}T_2 = \begin{bmatrix} 0 & -1 & 0 & 0 \\ 1 & 0 & 0 & 2 \\ 0 & 0 & 1 & 0 \\ 0 & 0 & 0 & 1 \end{bmatrix}, \quad
    {}^{2}T_3 = \begin{bmatrix} 0 & 1 & 0 & 0 \\ -1 & 0 & 0 & -1 \\ 0 & 0 & 1 & 0 \\ 0 & 0 & 0 & 1 \end{bmatrix}
\end{equation*}
\textbf{b) [5 pts]} ${}^{0}T_3 = {}^{0}T_1 {}^{1}T_2 {}^{2}T_3$.\\
\textbf{c) [13 pts]} ${}^{0}T_2 = \begin{bmatrix} 0&-1&0&0 \\ 0&0&-1&0 \\ 1&0&0&3 \\ 0&0&0&1 \end{bmatrix}$. ${}^{0}T_3 = {}^{0}T_2 {}^{2}T_3 = \begin{bmatrix} 1&0&0&1 \\ 0&0&-1&0 \\ 0&1&0&3 \\ 0&0&0&1 \end{bmatrix}$. (\textit{5 pts matriz intermedia, 5 pts orientación, 3 pts posición final}). \\
\textbf{d) [6 pts]} ${}^{0}p_3 = [1, 0, 3]^T$ (Posición del marco 3 en base 0). ${}^{0}R_3 = \begin{bmatrix} 1&0&0 \\ 0&0&-1 \\ 0&1&0 \end{bmatrix}$ (Orientación de los ejes del marco 3 vistos desde 0). \\
\textbf{e) [6 pts]} Verificación: $R^TR = I$; la fila base es homogénea $[0,0,0,1]$; o alcance físico $d_1+a_2+a_3$ es congruente. (\textit{Full credit por una estrategia correcta}).

\section*{Problema 4 (10 pts)[cite: 7]}
\textbf{a) [5 pts]} Físicamente imposible. El alcance máximo es $a_1+a_2 = 3$. La magnitud del vector dado es $\|p\| = 4$, lo cual excede el límite geométrico. \\
\textbf{b) [5 pts]} No es ortonormal. $R^TR \neq I$, o su determinante no es 1 (la norma de la segunda columna es 2).

\section*{Código de Verificación Docente (No incluir en examen)[cite: 7]}
\begin{lstlisting}
import numpy as np

def dh_standard(theta, d, a, alpha):
    ct, st = np.cos(theta), np.sin(theta)
    ca, sa = np.cos(alpha), np.sin(alpha)
    return np.array([
        [ct, -st*ca,  st*sa, a*ct],
        [st,  ct*ca, -ct*sa, a*st],
        [0.0,    sa,     ca,    d],
        [0.0,   0.0,    0.0,  1.0]
    ])

T01 = dh_standard(0.0, 1.0, 0.0, np.pi/2)
T12 = dh_standard(np.pi/2, 0.0, 2.0, 0.0)
T23 = dh_standard(-np.pi/2, 0.0, 1.0, 0.0)

T03 = T01 @ T12 @ T23
print(np.round(T03, 6))
\end{lstlisting}

\end{document}
